\documentclass[11pt]{article}
\usepackage[margin=1in]{geometry}
\usepackage{amsmath,amssymb,amsthm,booktabs,hyperref}
\newtheorem{lemma}{Lemma}
\newtheorem{proposition}[lemma]{Proposition}
\newtheorem{theorem}[lemma]{Theorem}
\newcommand{\R}{\mathbb R}
\newcommand{\dd}{\,d}
\newcommand{\Ea}{E_{\mathrm{arguments}}}
\newcommand{\Eh}{E_{\mathrm{horizontal}}}
\newcommand{\Eg}{E_{\mathrm{gaps}}}
\title{An ordinate-only reduction for the smoothed triple statistic}
\author{Work Package B extension: proof draft}
\date{October 3, 2026}
\begin{document}
\maketitle
\begin{abstract}
We give an exact transfer identity from the existing weighted full-zero
theorem, and prove that its two real-test kernel errors vanish.
The remaining complex-argument error is isolated, with a sufficient
positive-moment bound. Its vanishing is open: this draft does not prove
an ordinate-only correlation asymptotic.
\end{abstract}

\section{Conventions and the target}
Task: proof and source review. Assumptions: \texttt{UNCONDITIONAL}, and
\texttt{SUPPORT}\,$(h\leq1-\kappa,\ 0<\kappa<1)$ wherever the weighted
theorem or a Fourier test is used. Let $I$ index all occurrences of
nontrivial zeros $\rho_i=1/2+\delta_i+i\gamma_i$, with both ordinate
signs and full multiplicities. Every sum below is over all ordered
triples in $I^3$. Equal indices, equal complex zeros at different indices,
and distinct zeros with equal ordinates are all included.
Fix nonnegative $\omega\in C_c^\infty((1,2))$ of integral one, extended
by zero. For $T\geq2\pi e$ put
\[
 q=\log T,\qquad b=\log(T/(2\pi)),\qquad \ell=\frac b{2\pi},
 \qquad c_T=\frac{16}{3Tq}.
\]
Fix $\phi\in C_c^\infty(\R^2)$ supported in
$h(\xi,\eta)=\max(|\xi|,|\eta|,|\xi+\eta|)\leq s=1-\kappa<1$.
Use $\widehat F(\xi,\eta)=\int F(u,v)e^{-2\pi i(u\xi+v\eta)}\dd u\dd v$
and the entire extension
\[
 F(z_1,z_2)=\int_{\R^2}\phi(\xi,\eta)
              e^{2\pi i(\xi z_1+\eta z_2)}\dd\xi\dd\eta.
\]
For a triple, abbreviate
\[
 a=\gamma_{i_2}-\gamma_{i_1},\quad d=\gamma_{i_3}-\gamma_{i_1},\quad
 (u,v)=\ell(a,d),\quad r=\max(|u|,|v|),
\]
\[
 z=(u-i\ell(\delta_{i_2}+\delta_{i_1}),
       v-i\ell(\delta_{i_3}+\delta_{i_1})),\qquad
 w=\omega(\gamma_{i_1}/T).
\]
The conventional all-ordered ordinate statistic is
\begin{equation}\label{eq:ordinate}
 O_T[F]=\frac1{T\ell}\sum_{I^3}wF(u,v).
\end{equation}
Only the anchor is localized. No simplicity, ordinate-distinctness,
sharp cutoff, or exclusion of partial diagonals is imposed.

Write $J_\delta=J(a,d;\delta_{i_2},\delta_{i_3},\delta_{i_1})$,
$J_0=J(a,d;0,0,0)$, and $J_*=3\pi/8$. Explicitly, with
$c_1=a-i\delta_{i_2}$, $c_2=d-i\delta_{i_3}$, $c_3=i\delta_{i_1}$,
and $d_{mn}=c_m-c_n$,
\begin{equation}\label{eq:profile}
 J_\delta=\pi\frac{24+d_{12}^2+d_{13}^2+d_{23}^2}
 {(4+d_{12}^2)(4+d_{13}^2)(4+d_{23}^2)}
 =\int_\R\prod_{m=1}^3\frac1{1+(t-c_m)^2}\dd t.
\end{equation}
The profile and its integral representation are
\texttt{TRIPLE-KERNEL-PROFILE-001}. The baseline theorem
\texttt{TRIPLE-SMOOTHED-CORRELATION-001} states that
$R_T^J=c_T\sum wJ_\delta F(z)=S_3[F]+o(1)$ for fixed data.
Here, with $s_0(t)=\sin(\pi t)/(\pi t)$ and $s_0(0)=1$,
\begin{align*}
 R_2(t)&=1-s_0(t)^2,\\
 R_3(u,v)&=\det\begin{pmatrix}1&s_0(u)&s_0(v)\\s_0(u)&1&s_0(u-v)\\
 s_0(v)&s_0(u-v)&1\end{pmatrix},\\
 S_3[F]&=\int_{\R^2}FR_3\dd u\dd v+
 \int_\R R_2(t)\{F(0,t)+F(t,0)+F(t,t)\}\dd t+F(0,0).
\end{align*}
Pairings are bilinear; $F$ need not be real or symmetric.
\paragraph{Open target (\texttt{ORDINATE-TRIPLE-001}).}\label{target:ordinate}
Prove $O_T[F]=S_3[F]+o(1)$ for the fixed tests above.
The target has ledger status \texttt{idea}, not \texttt{proved-draft}.

\section{Exact transfer and convergence}
\begin{lemma}[Exact reduction; \texttt{ORDINATE-TRANSFER-001}]
\label{lem:transfer}
Define signed errors, without taking absolute values inside the sums, by
\begin{align}
 \Ea&=c_T\sum wJ_\delta\{F(z)-F(u,v)\},\label{eq:arguments}\\
 \Eh&=c_T\sum w(J_\delta-J_0)F(u,v),\label{eq:horizontal}\\
 \Eg&=c_T\sum w(J_0-J_*)F(u,v).\label{eq:gaps}
\end{align}
All these sums converge absolutely at each fixed $T$. Exactly,
\begin{align}
 R_T^J&=\frac bq O_T+\Ea+\Eh+\Eg,\label{eq:transfer}\\
 O_T-S_3[F]&=\frac qb\{R_T^J-S_3[F]-\Ea-\Eh-\Eg\}
                   +\left(\frac qb-1\right)S_3[F].\label{eq:main-transfer}
\end{align}
Zero cutoffs may tend independently to infinity at fixed $T$ and fixed
test; only afterwards is $T$ allowed to grow.
\end{lemma}
\begin{proof}
For real $x$ and $|e|\leq1/2$,
$|1+(x-ie)^2|\geq1+x^2-e^2\geq\tfrac34(1+x^2)$.
The first two $e$ derivatives of its reciprocal are bounded by
$C(1+x^2)^{-1}$: differentiate the rational function and use this
inequality, with $|x-ie|\ll(1+x^2)^{1/2}$.
Consequently $J$ and its horizontal derivatives through order two are
uniformly bounded for all real gaps and the closed parameter cube
$[-1/2,1/2]^3$. Indeed, bound two real kernels by one and integrate the
third. Differentiation under the integral is justified locally in the
parameters by an integrable dominating function; the resulting bounds
are global. Colliding centers cause no exception.

The anchor has finitely many occurrences. The unit-interval count
\texttt{PAIR-COUNT-001} gives $O(\log(|t|+2))$ zeros in a unit interval
near $t$, with multiplicity. At fixed $T$ the imaginary coordinates
of $z$ are bounded by $\ell$. Integrating the compact smooth Fourier
integral by parts gives, for every $N$, a bound
$|F(z)|+|F(u,v)|\leq C_{T,N,\phi}(1+r)^{-N}$.
Two unit-interval summations therefore converge absolutely. This
also justifies reindexing by horizontal reflection and all cutoff
limits here. Telescope each summand. Finally
$c_TJ_*=2\pi/(Tq)=(b/q)/(T\ell)$, proving both identities.
\end{proof}

\section{Two unconditional counting inputs}
These are imported source statements, not numerical project certificates.
Their upstream analytic proofs have not been reconstructed in this extension.

\begin{proposition}[Close-pair input; \texttt{ORDINATE-PAIR-INPUT-001}]
\label{prop:pair-input}
For $V$ sufficiently large, $L_V=\log V/(2\pi)$ and
$1\leq\lambda\leq\sqrt{L_V}$,
\begin{equation}\label{eq:pair-input}
 \sum_{\substack{0<\gamma,\gamma'\leq V\\
                  |\gamma-\gamma'|L_V\leq\lambda}}
 \left(1-\frac{|\gamma-\gamma'|L_V}{\lambda}\right)
 \ll V\log V\,\lambda.
\end{equation}
The sum includes all ordered occurrences, including equal ordinates.
\end{proposition}
\begin{proof}[Source and deduction]
GLSS~\cite{GLSS} (\textbf{PREPRINT}), equation (6.1), gives the left side as
$VL_V\{\lambda+\log(2+\lambda)/(\pi^2\lambda)
+O(\sqrt{\log(2+\lambda)}/\lambda)\}$ uniformly for
$0<\lambda^2\leq L_V$. For $\lambda\geq1$ each displayed term is
$O(\lambda)$. Their Section 6 derives this from the unconditional
Propositions 1--2; their PCC-dependent main theorem and the RH improvement
in Remark 2 are not used. No explicit constant is asserted here.
\end{proof}

\begin{proposition}[Density input; \texttt{ORDINATE-DENSITY-INPUT-001}]
\label{prop:density-input}
Let $N(\sigma,V)$ count $\beta>\sigma$, $0<\gamma\leq V$, with
multiplicity. Uniformly for $0\leq D\leq0.331$ and
$V\geq H_0=3.0610046\cdot10^{10}$,
\begin{equation}\label{eq:density-input}
 N(1/2+D,2V)-N(1/2+D,V)
 \ll V^{1-D/4}\log V+(\log V)^2.
\end{equation}
\end{proposition}
\begin{proof}[Source and deduction]
Simoni\v c, Theorem 1, consulted in arXiv:1910.08274v2, p.~2
(published in J.~Math.~Anal.~Appl.~491 (2020), 124303), bounds this
difference by
\[
10395.2 V^{1-D/4}\log V+1.104\log^2V+
0.173\log V\log\log V+0.51\log V.
\]
This implies \eqref{eq:density-input} with an effective absolute constant.
The consulted author text is distinguished from the journal version.
Its external finite-height and numerical inputs are inherited, not replayed.
No zero in our sums is replaced by a critical-line zero.
\end{proof}

\begin{lemma}[Triple counts; \texttt{ORDINATE-COUNT-001}]
\label{lem:counts}
Let $H_T(R)$ count all triples with $T\leq\gamma_{i_1}\leq2T$
and $r\leq R$. For sufficiently large $T$,
\begin{align}
 H_T(R)&\ll Tq^2R^2 &&(1\leq R\leq R_0=q^{1/4}),\label{eq:small-count}\\
 H_T(R)&\ll Tq^3(1+R)^4 &&(R\geq1).\label{eq:global-count}
\end{align}
Put $A_{20}(F)=\sup_{u,v}(1+\max(|u|,|v|))^{20}|F(u,v)|$ and
$W=\|\omega\|_\infty$. Then
\begin{align}
 \frac1{Tq}\sum_{r\leq R_0}w|F(u,v)|&\ll WA_{20}(F)q,\label{eq:mass}\\
 \frac1{Tq}\sum_{r>R_0}w|F(u,v)|&\ll WA_{20}(F)q^{-2}.\label{eq:tail}
\end{align}
\end{lemma}
\begin{proof}
For the small range $R/\ell<1$. All participating zeros have positive
ordinates at most $3T$, and each anchor has $O(q)$ possible partners.
Apply \eqref{eq:pair-input} at $V=3T$ and
$\lambda=2RL_{3T}/\ell$. For large $T$ this belongs to
$[1,\sqrt{L_{3T}}]$ throughout $R\leq R_0$. Every pair at distance
at most $R/\ell$ has triangular weight at least $1/2$.
Thus there are $O(TqR)$ anchor--partner pairs. Multiplying by the
$O(q)$ choices for the third occurrence proves even $O(Tq^2R)$,
and hence \eqref{eq:small-count}.

Globally there are $O(Tq)$ anchors and at most
$C(1+R/\ell)\log(T+R/\ell+2)$ partners per anchor, by unit counts
and conjugation for negative ordinates. Since $\ell\geq1$ eventually,
this is at most $Cq(1+R)^2$, giving \eqref{eq:global-count}.
For \eqref{eq:mass}, use $r\leq1$ and dyadic shells up to $R_0$;
their normalized bounds are $CWA_{20}q\sum_{k\geq0}2^{(2-20)k}$.
A shell cut by $R_0$ uses $H_T(R_0)$ and has the same bound.
For the tail use shells $2^kR_0<r\leq2^{k+1}R_0$ and
\eqref{eq:global-count}; the result is
$CWA_{20}q^2R_0^{4-20}=CWA_{20}q^{-2}$.
\end{proof}

\section{Removing the real-test kernel}
\begin{theorem}[Real-test errors; \texttt{ORDINATE-REAL-KERNEL-001}]
\label{thm:real-kernel}
For every Schwartz $F$ on $\R^2$ and the smoothing above, as $T\to\infty$,
\begin{equation}\label{eq:kernel-errors}
 |\Eg|\ll \frac{WA_{20}(F)}q,\qquad
 |\Eh|\ll WA_{20}(F)\frac{(\log q)^2}q.
\end{equation}
Constants and the lower threshold for $T$ are independent of $F,\omega$.
No Fourier support is needed for these two estimates.
\end{theorem}
\begin{proof}
The rational profile is smooth near $(a,d)=(0,0)$, is even under
$(a,d)\mapsto(-a,-d)$, and equals $J_*$ at the origin. Together
with its global bound this gives
$|J_0-J_*|\leq C\min(a^2+d^2,1)$.
On $r\leq R_0$ use $C r^2/\ell^2$ and the same dyadic argument as
\eqref{eq:mass}: the bound is
\[
CWA_{20}q\ell^{-2}\sum_{k\geq0}2^{(4-20)k}\ll WA_{20}/q.
\]
On the tail use \eqref{eq:tail}.

For the horizontal error, reflect all three indices simultaneously by
$\rho\mapsto1-\bar\rho$. This preserves ordinates, multiplicities,
$w$, $F(u,v)$ and every index coincidence pattern, while sending
$\delta\mapsto-\delta$. Hence in \eqref{eq:horizontal} one may
replace $J_\delta-J_0$ by
$\tfrac12(J_\delta+J_{-\delta})-J_0$.
Taylor's formula and the uniform second-derivative bounds in the proof
of Lemma~\ref{lem:transfer} give
\[
 \left|\tfrac12(J_\delta+J_{-\delta})-J_0\right|
 \leq C\max_j|\delta_{i_j}|^2.
\]
This is a reindexing of the sum, not a termwise positivity assertion.

Take $D=64\log q/q$, eventually in $(0,0.331)$.
On $r\leq R_0$ with every $|\delta_{i_j}|\leq D$, equation
\eqref{eq:mass} bounds the error by $CWA_{20}qD^2$.
For the remaining tuples some occurrence with $|\delta|>D$ lies in
$[T-1,2T+1]$. Reflection and \eqref{eq:density-input}, applied to the
three windows $(T/2,T],(T,2T],(2T,4T]$, show that their number is
\[
 B_T(D)\ll Tq\exp\{-D\log(T/2)/4\}+q^2
          \ll Tq^{-7}+q^2.
\]
The last inequality uses $\log(T/2)/q\geq1/2$ and thus an exponent
at most $-8\log q$. Each such occurrence belongs to at most $Cq^2$
of our near tuples: fix its slot, and use two unit-interval counts
within distance two. Union over the three slots only increases this
bound. Boundedness of the kernel and $|F|\leq A_{20}$ give a normalized
contribution $CWA_{20}(q^{-6}+q^3/T)$.
The tail is $CWA_{20}q^{-2}$ by \eqref{eq:tail}.
Combining these proves the second bound. Strict $|\delta|>D$ matches
the source's strict $\beta>\sigma$; equality belongs to the good set.
\end{proof}

\section{What remains at the complex arguments}
Define the nonnegative quantity (not assumed small)
\begin{equation}\label{eq:moment}
 M_T(s)=\frac1{Tq}\sum_{I^3}w(1+r)^{-20}
 x^2(1+x)^{22}e^{2sx},\qquad x=q\max_j|\delta_{i_j}|,
\end{equation}
and $P_{22}(\phi)=\sum_{|\alpha|\leq22}\|\partial^\alpha\phi\|_1$.
\begin{lemma}[Reflected argument bound; \texttt{ORDINATE-ARGUMENT-BOUND-001}]
\label{lem:arguments}
Uniformly for the fixed-support class above and $T$ sufficiently large,
\begin{equation}\label{eq:argument-bound}
 |\Ea|\leq C P_{22}(\phi)M_T(s).
\end{equation}
The moment is finite for every fixed $T$. Its convergence to zero has
not been established.
\end{lemma}
\begin{proof}
For one triple set $y=(\ell(\delta_{i_2}+\delta_{i_1}),
\ell(\delta_{i_3}+\delta_{i_1}))$ and
$G(t)=J(a,d;t\delta)\{F((u,v)-ity)-F(u,v)\}$ for $-1\leq t\leq1$.
Simultaneous reflection replaces its contribution by
$(G(1)+G(-1))/2$. Since $G(0)=0$, Taylor's integral remainder bounds
this by $\tfrac12\sup_{|t|\leq1}|G''(t)|$.
Inside the Fourier integral the real exponent is $tA$, where
\[
 A=b\{\delta_{i_2}\xi+\delta_{i_3}\eta+
                   \delta_{i_1}(\xi+\eta)\}.
\]
The exact identity $|\xi|+|\eta|+|\xi+\eta|=2h(\xi,\eta)$ gives
$|A|\leq2sx$, since $b/q<1$. Also $|\nabla A|\leq Cx$.
Apply $(1-\Delta_{\xi,\eta})^{10}$ in the Fourier integral.
For $k=0,1,2$, the derivatives of $A^ke^{tA}\phi$ have $L^1$ norm
at most $CP_{22}(\phi)x^k(1+x)^{20}e^{2sx}$.
The integration-by-parts denominator is
$(1+4\pi^2(u^2+v^2))^{10}$, so these estimates supply decay
$(1+r)^{-20}$. Compact support makes the boundary terms zero.
In $G''$, the three product-rule terms involve respectively $J''$,
$2J'$ and $J$, where the first two derivatives along $t\delta$ are
$O((x/q)^2)$ and $O(x/q)$. Bound the undifferentiated difference
of tests by the sum of their absolute values. It follows that
\[
 |G''(t)|\leq CP_{22}(\phi)(1+r)^{-20}
                  x^2(1+x)^{22}e^{2sx}.
\]
Sum this estimate using Lemma~\ref{lem:transfer}. Finiteness of
\eqref{eq:moment} follows from $x\leq q/2$ at fixed $T$ and the
global count \eqref{eq:global-count}.
\end{proof}

\paragraph{Exact remaining criterion.}\label{target:arguments}
For each fixed admissible test and smoothing, \eqref{eq:main-transfer},
the baseline theorem, and Theorem~\ref{thm:real-kernel} imply
\[
 O_T[F]=S_3[F]+o(1)\quad\Longleftrightarrow\quad\Ea=o(1).
\]
This is a consequence of \texttt{ORDINATE-TRANSFER-001}, not a new
assertion that either limit holds. In particular $M_T(s)=o(1)$ would
be sufficient. The other named errors are the inherited
$R_T^J-S_3[F]$, the two bounds \eqref{eq:kernel-errors}, and the exact
scale error $(q/b-1)S_3[F]$; all enter with the signs in
\eqref{eq:main-transfer}. For varying tests one must make each of
these bounds, including $P_{22}(\phi)M_T(s)$, tend to zero.

\paragraph{Why the density input does not finish this step.}
Small unscaled $\delta$ does not give small $q\delta$.
The density bound gives no vanishing bound for the fraction at
$|\delta|=c/q$, with fixed $c>0$. The cutoff $64\log q/q$ used for
the kernel leaves shifts of size $O(\log q)$ in the test arguments.
It therefore cannot be used to make those arguments real by continuity.
Signed cancellation might still prove $\Ea=o(1)$ without the positive
moment condition; we assert no impossibility theorem.

\paragraph{Finite sensitivity check (\texttt{SYNTHETIC\_MODEL}).}
Put occurrences with $\delta=\pm c/q$ at one positive ordinate and
include their negative conjugates. At zero vertical gaps the kernel
tends to $J_*$ as $q\to\infty$. Averaging a Fourier mode over the
eight ordered choices of positive-ordinate signs gives the multiplier
\[
 \cosh((b/q)c\xi)\cosh((b/q)c\eta)
                       \cosh((b/q)c(\xi+\eta)).
\]
Its limit is greater than one when $c>0$ and $(\xi,\eta)\ne(0,0)$.
This proves pointwise sensitivity after symmetry averaging. A finite
quartet has vanishing mass after normalization by $Tq$; it is not a
counterexample to the zeta asymptotic. A Fourier mode is used only for
this finite identity, not as an admissible Schwartz test.

\section{Provenance and scope}
This extension uses \texttt{PAIR-COUNT-001}, the recorded kernel profile,
and the baseline weighted theorem. The two new source statements are
imported as unconditional inputs with their ranges intact. Neither
PCC, RH, simplicity nor an all-zero critical-line specialization occurs
in the reduction. The prior line-by-line audit and archived reproduction
cover the weighted baseline; they do not certify this extension.
The new exact regression script checks finite algebra and normalization,
not asymptotic analysis or the imported theorems.
All limits here fix $\omega,\phi,\kappa$ first, remove zero cutoffs by
absolute convergence at fixed $T$, and then send $T\to\infty$.
Internal Fourier axes and the origin are included; the outer support
boundary is excluded by $\kappa>0$. No new endpoint extension is used.

\begin{thebibliography}{9}
\bibitem{GLSS}
D.~A.~Goldston, J.~Lee, J.~Schettler, A.~I.~Suriajaya,
\emph{Pair Correlation Conjecture for the zeros of the Riemann
zeta-function I: simple and critical zeros},
\href{https://arxiv.org/abs/2503.15449v4}{arXiv:2503.15449v4},
PREPRINT, 2026. Section 6, equation (6.1).
\bibitem{Simonic}
A.~Simoni\v c, \emph{Explicit zero density estimate for the Riemann
zeta-function near the critical line}, J.~Math.~Anal.~Appl.~491 (2020),
124303, \href{https://doi.org/10.1016/j.jmaa.2020.124303}{DOI}.
Consulted \href{https://arxiv.org/abs/1910.08274v2}{arXiv:1910.08274v2},
Theorem 1, p.~2; proof Section 4.5.
\end{thebibliography}
\end{document}
